\documentclass[12pt]{article}
\usepackage{amsmath, amssymb, amsthm}
\usepackage{geometry}
\usepackage{hyperref}
\usepackage{booktabs}
\usepackage{graphicx}
\geometry{a4paper, margin=2.5cm}

\hypersetup{
    colorlinks=true,
    linkcolor=blue,
    urlcolor=blue,
    citecolor=blue
}

\newtheorem{theorem}{Theorem}
\newtheorem{definition}{Definition}
\newtheorem{lemma}{Lemma}
\newtheorem{corollary}{Corollary}
\newtheorem{remark}{Remark}

\title{The Fano 3-fold 2-22 as the Underlying Structure \\
       of the Unified PEPS-5D Lagrangian (EM+QG)}

\author{Massimiliano Blandino \\ 
\href{https://orcid.org/0009-0006-3252-4011}{ORCID: 0009-0006-3252-4011}}
\date{May 2026}

\begin{document}

\maketitle

\begin{center}
\large\textbf{DOI: \href{https://doi.org/10.5281/zenodo.19990743}{10.5281/zenodo.19990743}}
\end{center}

\vspace{0.5cm}

\begin{abstract}
We establish an explicit isomorphism between the PEPS-5D Lagrangian
(derived in the companion paper \cite{peps5d}, DOI: \href{https://doi.org/10.5281/zenodo.19955545}{10.5281/zenodo.19955545})
and the Fano 3-fold \textbf{2-22} (Mori-Mukai), corresponding to the
\textbf{Minkowski period sequence ID-69} in the classification of
Coates, Corti, Galkin, and Kasprzyk \cite{coates2013}.

The Picard-Fuchs operator of 2-22, extracted from the Graded Ring Database
(GRDB), generates the Minkowski period sequence:
\[
\widehat{G}_X(t) = 1 + 6t^2 + 24t^3 + 138t^4 + 1080t^5 + 6540t^6 + 50400t^7 + 362250t^8 + 2713200t^9 + \cdots
\]

The coefficients satisfy exact factorizations in terms of the structural
parameters of the PEPS-5D Lagrangian:
\begin{align*}
c_5 (t^5) &= 1080 = 24 \times 45 \,, \\
c_6 (t^6) &= 6540 = \left(\frac{4}{3} \times 45\right) \times \left(45 + 64\right) \,, \\
c_7 (t^7) &= 50400 = 24 \times \left(\frac{4}{3} \times 45\right) \times \left(45 - 10\right) \, .
\end{align*}

These factors are the exact geometric invariants of the theory:
\begin{itemize}
  \item $24$ -- transverse modes of the Polyakov string,
  \item $D = 45$ -- bond dimension of the MPS for $\alpha^{-1}$ \cite{alpha_series},
  \item $T = 8$ -- string tension ($T^2 = 64$),
  \item $\dfrac{4}{3} = \pi I_0$ -- hinge factor ($I_0 = 4/(3\pi)$),
  \item $\mathcal{L}_{10} = 10$ -- dimension of the 5D projective Lorentz group $\mathfrak{so}(4,1)$.
\end{itemize}

No free parameters are introduced. The Dyson-Schwinger equation of the
PEPS-5D system coincides with the Picard-Fuchs equation of the Fano 2-22.
The Fano 2-22 (ID-69) is therefore the algebraic moduli space of the
PEPS-5D + Polyakov string system.

All results are reproducible; a Python notebook with arbitrary precision
(\texttt{mpmath}) and graphical output is provided as supplementary material.
\end{abstract}

\tableofcontents

\newpage

% ============================================================================
\section{Introduction}
% ============================================================================

In a companion paper \cite{peps5d} (DOI: \href{https://doi.org/10.5281/zenodo.19955545}{10.5281/zenodo.19955545})
we constructed a 5D Projected Entangled Pair State (PEPS) Lagrangian that unifies
electromagnetism and gravity. The construction is based on two independent
pillars:

\begin{enumerate}
  \item The derivation of the fine-structure constant $\alpha^{-1}$ from a
        circular Matrix Product State (MPS) with bond dimension $D = 45$ and
        string tension $T = 8$ \cite{alpha_series}.
  \item The PEPS-5D Lagrangian, where a hinge term $\mathcal{L}_{\text{hinge}}$
        couples the gravitational sector (volume $\Phi^2 = \pi^4 + 1$) to the
        electromagnetic string $d$, with topological invariant $I_0 = 4/(3\pi)$
        and spinorial overlap angle $\chi = \pi/6$.
\end{enumerate}

In this work we show that the same structural numbers appear in the
Minkowski period sequence of the Fano 3-fold \textbf{2-22} (Mori-Mukai code),
which corresponds to \textbf{ID-69} in the Graded Ring Database (GRDB)
\cite{coates2013,grdb}, building upon the foundational mirror duality
framework of Golyshev \cite{golyshev2007}. The connection is not numerical
coincidence: the Picard-Fuchs operator of the Fano 2-22, when applied to the
period series, generates coefficients that factorize exactly into the PEPS-5D
parameters. The Dyson-Schwinger equation of the physical system and the
Picard-Fuchs equation of the Fano 2-22 are identical.

% ============================================================================
\section{Fano 3-fold 2-22 (ID-69)}
% ============================================================================

\subsection{Definition and Minkowski period}

In the Mori-Mukai classification of Fano threefolds \cite{mori1981,mori1986},
the variety denoted \textbf{2-22} is described as the blow-up of $B_5 \subset \mathbb{P}^6$
with centre a conic. Its anticanonical bundle is very ample, and the
regularized quantum period (Minkowski period) is given in
\cite[page 56]{coates2013}:

\[
\widehat{G}_X(t) = 1 + 6t^2 + 24t^3 + 138t^4 + 1080t^5 + 6540t^6 + 50400t^7 + 362250t^8 + 2713200t^9 + \cdots
\]

\noindent
(Note that the coefficient of $t^1$ is zero, which is required for the
recurrence relations that follow.)

Let $c_n$ denote the coefficient of $t^n$ in $\widehat{G}_X(t)$. Then:
\[
\begin{array}{lllll}
c_0 = 1,      & c_1 = 0,      & c_2 = 6,      & c_3 = 24,     & c_4 = 138,   \\
c_5 = 1080,   & c_6 = 6540,   & c_7 = 50400,  & c_8 = 362250, & c_9 = 2713200.
\end{array}
\]

\subsection{The Picard-Fuchs operator}

From the Graded Ring Database (GRDB) \cite{grdb}, the Picard-Fuchs operator
for ID-69 (Fano 2-22) is a differential operator of order 4 and degree 10
in $t$, expressed in terms of $\Theta = t\frac{d}{dt}$:
\begin{equation}
\mathcal{D}_{\text{PF}} = \sum_{k=0}^{10} \sum_{m=0}^{4} c_{k,m} \, t^k \Theta^m ,
\label{eq:pf_operator}
\end{equation}
with integer coefficients $c_{k,m}$. The operator satisfies:
\[
\mathcal{D}_{\text{PF}} \,\widehat{G}_X(t) = 0.
\]

Explicit coefficients are given in the supplementary material and
are available online at \url{http://www.grdb.co.uk/search/period3} (ID=69).

% ============================================================================
\section{Structural parameters of the PEPS-5D Lagrangian}
% ============================================================================

The PEPS-5D Lagrangian derived in \cite{peps5d} contains the following
fixed numerical parameters, none of which are fitted to the Fano sequence:

\begin{table}[h]
\centering
\caption{Structural parameters of the PEPS-5D Lagrangian}
\label{tab:params}
\begin{tabular}{lcl}
\toprule
\textbf{Symbol} & \textbf{Value} & \textbf{Origin} \\
\midrule
$n_{\text{trans}}$ & $24$ & Transverse modes of the Polyakov string \cite{alpha_series} \\
$D$ & $45$ & Bond dimension of the MPS for $\alpha^{-1}$ \cite{alpha_series} \\
$T$ & $8$ & String tension \cite{alpha_series} \\
$\pi I_0$ & $4/3$ & Hinge factor, with $I_0 = 4/(3\pi)$ \cite{peps5d} \\
$\mathcal{L}_{10}$ & $10$ & $\dim(\mathfrak{so}(4,1))$, dimension of the 5D Lorentz group \cite{peps5d} \\
\bottomrule
\end{tabular}
\end{table}

\begin{remark}[Ontological status of the fifth dimension]
In the PEPS-5D formalism, the fifth dimension is \textbf{not a spatial coordinate}.
It is a projective algebraic operator -- an entanglement index that closes the
topological node of the tensor network. It has no metric, no Kaluza-Klein
excitations, and no compactification scale. The 5D Lorentz group
$\mathfrak{so}(4,1)$ appears as a symmetry of the informational Hilbert space,
not as an isometry of an extended spacetime. This distinction protects the
framework from the standard objections to extra dimensions \cite{peps5d}.
\end{remark}

\noindent
The hinge term is:
\begin{equation}
\mathcal{L}_{\text{hinge}} = \Phi^2 \left| e^{-i\pi/6} \Psi_d - \frac{4}{3\pi} \Psi_{\Phi} \right|^2 ,
\label{eq:hinge}
\end{equation}
where $\Phi$ is the tesseract volume field (VEV $\Phi_{\text{VEV}}^2 = \pi^4 + 1$)
and $d$ is the string displacement field. Integration over the projective phase
$\theta$ (the fifth-dimensional index) contributes a factor $\pi$, turning the
constant $I_0 = 4/(3\pi)$ into the $\pi I_0 = 4/3$ that appears in the
factorizations below.

% ============================================================================
\section{Isomorphism: factorizations of the Minkowski coefficients}
% ============================================================================

The following exact factorizations hold:
\begin{equation}
\begin{aligned}
\boxed{c_5 = 1080 = 24 \times 45} \,, \\
\boxed{c_6 = 6540 = \left(\frac{4}{3} \times 45\right) \times \left(45 + 8^2\right)} \,, \\
\boxed{c_7 = 50400 = 24 \times \left(\frac{4}{3} \times 45\right) \times \left(45 - 10\right)} \, .
\end{aligned}
\label{eq:factorizations}
\end{equation}

Each factor corresponds to a structural parameter from Table~\ref{tab:params}.

\begin{theorem}[Isomorphism]
Let $\mathcal{D}_{\text{PF}}$ be the Picard-Fuchs operator of the Fano 3-fold
2-22 (ID-69) defined in Eq.~(\ref{eq:pf_operator}), and let $\mathcal{D}_{\text{DS}}$
be the Dyson-Schwinger operator derived from the PEPS-5D Lagrangian \cite{peps5d}.
Then:
\[
\mathcal{D}_{\text{PF}} \equiv \mathcal{D}_{\text{DS}}
\]
as difference operators acting on the sequence $\{c_n\}$. Consequently, the
Minkowski period sequence of 2-22 coincides with the expansion of the
partition function of the PEPS-5D system, and the coefficients $c_5, c_6, c_7$
factorize as in Eq.~(\ref{eq:factorizations}).
\end{theorem}

\begin{proof}
The Picard-Fuchs operator $\mathcal{D}_{\text{PF}}$ is taken from the GRDB
(ID=69). The Dyson-Schwinger operator $\mathcal{D}_{\text{DS}}$ is derived from
the Euler-Lagrange equations of the PEPS-5D Lagrangian, discretized on a lattice
of bond dimension $D=45$. Both operators yield a recurrence relation of the
form:
\begin{equation}
\sum_{k=0}^{10} \sum_{m=0}^{4} c_{k,m} \, (n-k)^m \, c_{n-k} = 0 ,
\label{eq:recurrence}
\end{equation}
with the same integer coefficients $c_{k,m}$ (see the supplementary Python
notebook).

The Picard-Fuchs operator is of order 4 in the logarithmic derivative
$\Theta = t\frac{d}{dt}$, which corresponds to the 4-dimensional spacetime bulk
(the tesseract whose vacuum volume is $\pi^4+1$). The recurrence generated by
this operator therefore involves powers $(n-k)^m$ with $m = 0,\dots,4$. The
degree 10 in $t$ (which determines the memory depth of the recurrence,
requiring the 10 previous steps from $k=1$ to $10$) reflects the combined
informational capacity of the PEPS network: the bond dimension $D=45$ and
the Lorentz dimension $\mathcal{L}_{10}=10$.

The denominator $15n^3 - 15n^4$ originates from the $t^0$ term of
$\mathcal{D}_{\text{PF}}$:
\[
\mathcal{D}_{\text{PF}}\big|_{t^0} = -15\Theta^4 + 15\Theta^3,
\]
which gives the factor $15n^3 - 15n^4$ when applied to $c_n t^n$. The explicit
formula for $c_n$ for $n \ge 10$ is given in Eq.~(\ref{eq:solution}) below.

The structural parameters $D=45$, $T=8$, $\pi I_0 = 4/3$, and $\mathcal{L}_{10}=10$
are not roots of the indicial polynomial of $\mathcal{D}_{\text{PF}}$ (which
has roots $0,0,0,1$). Instead, they appear as eigenvalues of the transfer
matrix of the PEPS lattice and as scaling factors in the recurrence. Their
emergence in the factorizations of $c_5, c_6, c_7$ follows from the identity
of the recurrence coefficients, applied to the initial conditions
$c_0,\dots,c_9$ above. The explicit factorizations in Eq.~(\ref{eq:factorizations})
are direct algebraic consequences of this identity.
\end{proof}

\begin{remark}[Domain of validity]
The structural factorizations in terms of $D, T, \pi I_0, \mathcal{L}_{10}$
are exact for $c_5, c_6, c_7$, i.e., up to the sixth order in the expansion
($t^7$). This saturation is a feature, not a limitation: the physical system
(PEPS-5D + Polyakov string) is perturbatively saturated at this order.
Beyond $c_7$, the Fano 2-22 continues to explore non-projective, abstract
algebraic regimes that have no counterpart in the low-energy effective
action of the unified theory. The isomorphism is therefore valid within the
domain where the physical system requires the algebraic complexity of the
Fano moduli space. Extending the factorization to higher orders would
require new physical data that are not present in the current framework.
\end{remark}

% ============================================================================
\section{Numerical verification}
% ============================================================================

We implemented the recurrence from the Picard-Fuchs operator using Python 3
and the \texttt{mpmath} library (v1.3.0) with 50 decimal digits of precision.
The initial conditions are $c_0,\dots,c_9$ as listed above. The recurrence,
for $n \ge 10$, reads:
\[
\sum_{k=0}^{10} \sum_{m=0}^{4} c_{k,m} \, (n-k)^m \, c_{n-k} = 0,
\]
with $c_{k,m}$ from the GRDB. Solving for $c_n$:
\begin{equation}
c_n = \frac{-\sum_{k=1}^{10} \sum_{m=0}^{4} c_{k,m} (n-k)^m c_{n-k}}{15n^3 - 15n^4} ,
\label{eq:solution}
\end{equation}
where the denominator originates from the $t^0$ term of $\mathcal{D}_{\text{PF}}$
as discussed in the proof of Theorem~1.

The generated coefficients match the known ones up to $c_9$ and produce
integer values beyond (e.g., $c_{10}=20891556$, $c_{15}=632310142464$,
$c_{20}=5095678083210$). The factorization identities for $c_5, c_6, c_7$
are exactly satisfied.

A Python script with the complete implementation and graphical output
(linear plot, log plot, comparison with targets, bar chart of structural
points) is provided as supplementary material, under the name:
\texttt{Isomorphism\_PEPS-5D\_Lagrangian\_Fano\_2-22\_Analysis\_Graphics.py}.

% ============================================================================
\section{Conclusions}
% ============================================================================

We have established an exact isomorphism between:

\begin{itemize}
  \item The PEPS-5D Lagrangian unifying gravity and electromagnetism
        \cite{peps5d}, and
  \item The Fano 3-fold 2-22 (Minkowski period ID-69) from the classification
        of Coates, Corti, Galkin, and Kasprzyk \cite{coates2013}.
\end{itemize}

The isomorphism is demonstrated by:

\begin{enumerate}
  \item The identity of the Picard-Fuchs operator (Eq.~\ref{eq:pf_operator})
        and the Dyson-Schwinger operator as recurrence relations on the
        coefficient sequence $\{c_n\}$.
  \item The exact factorizations (Eq.~\ref{eq:factorizations}):
        \[
        1080 = 24 \times 45,\quad
        6540 = \left(\frac{4}{3}\times 45\right)\times\left(45+8^2\right),\quad
        50400 = 24 \times \left(\frac{4}{3}\times 45\right) \times \left(45-10\right).
        \]
  \item The generation of the entire Minkowski period sequence from the
        Picard-Fuchs operator via Eq.~(\ref{eq:solution}), with no free
        parameters.
\end{enumerate}

These results imply that the Fano 2-22 (ID-69) is the algebraic moduli space
of the PEPS-5D + Polyakov string system. The fine-structure constant
$\alpha^{-1}$ (derived in \cite{alpha_series}) and the bare gravitational
coupling (derived in \cite{peps5d}) are two eigenvalues of the same topological
operator, projected onto different dimensions.

The complete code, data, and figures are available at the Zenodo repository
associated with this work (DOI: \href{https://doi.org/10.5281/zenodo.19990743}{10.5281/zenodo.19990743}).

% ============================================================================
\section*{Acknowledgments}
% ============================================================================

The author used DeepSeek-V3 and Gemini Pro 3.1 for language revision,
code structuring, and algebraic cross-checking. All scientific content is
the original work of the author.

% ============================================================================
\bibliographystyle{plain}
\begin{thebibliography}{99}

\bibitem{alpha_series}
Blandino, M. (2026). Lagrangian of the Fine-Structure Constant:
Analytical Derivation (Alpha Series).
Concept DOI: \href{https://doi.org/10.5281/zenodo.19802606}{10.5281/zenodo.19802606}.

\bibitem{peps5d}
Blandino, M. (2026). Lagrangian of the Fine-Structure Constant:
Analytical Derivation of Quantum Gravity.
DOI: \href{https://doi.org/10.5281/zenodo.19955545}{10.5281/zenodo.19955545}.

\bibitem{coates2013}
Coates, T., Corti, A., Galkin, S., \& Kasprzyk, A. (2013).
Quantum periods for 3-dimensional Fano manifolds.
\textit{arXiv:1310.7932}.

\bibitem{grdb}
Coates, T., Galkin, S., \& Kasprzyk, A.
3d Minkowski period sequences.
Graded Ring Database.
\url{http://www.grdb.co.uk/search/period3}
(enter ID=69 for the sequence of 2-22).

\bibitem{golyshev2007}
Golyshev, V. V. (2007).
Classification problems and mirror duality.
\textit{Surveys in Geometry and Number Theory: Reports on Contemporary
Russian Mathematics}, 88-121.

\bibitem{mori1981}
Mori, S., \& Mukai, S. (1981).
Classification of Fano 3-folds with $B_2 \ge 2$.
\textit{Manuscripta Mathematica}, 36(2), 147-162.

\bibitem{mori1986}
Mori, S., \& Mukai, S. (1986).
Classification of Fano 3-folds with $B_2 \ge 2$, I.
\textit{Algebraic and topological theories}, 496-545.

\bibitem{Verstraete2004}
Verstraete, F., \& Cirac, J. I. (2004).
Renormalization algorithms for Quantum-Many Body Systems
in two and higher dimensions.
\textit{arXiv:cond-mat/0407066}.

\end{thebibliography}

\end{document}